\section{Related Work}
\label{sec:related}

\subsection{Graph Neural Networks on Relational Databases}
RDL represents a relational database as a heterogeneous temporal graph with
one node per row and one edge per PK--FK reference, and learns to aggregate
information across tables with GNNs \citep{fey2024rdl}; RelBench
standardizes databases, temporal prediction tasks, and evaluation splits for
this setting \citep{robinson2024relbench}. RelGNN introduces atomic routes
and composite message passing to handle many-to-many relationships
\citep{chen2025relgnn}, and RelGT tokenizes each node into features, type,
hop distance, time, and local structure for a graph transformer
\citep{dwivedi2026relgt}. In RDL and RelGNN, messages flow only along PK--FK
references, so the database structure rather than learned attention decides
which rows interact. These models, however, encode the structure through
parameters tied to one schema: column encoders for each table, node and edge
types defined by the schema, separate weights for each relation or atomic
route, and an output head for each task
\citep{fey2024rdl,robinson2024relbench,chen2025relgnn}. They are therefore
trained anew for every database and task. \model{} keeps their
structure-following aggregation but replaces these schema-specific parameters
with counterparts shared across databases.

\subsection{Relational Foundation Models}
Relational foundation models reuse one pretrained model across databases.
Griffin pretrains shared encoders and decoders with relational message
passing, gathering the cells of each row through cross-attention over column
names, and is fine-tuned on each downstream task \citep{wang2025griffin}.
\RT{} removes fine-tuning: each cell is a token that carries an embedding of
its column and table names, the task is specified through the task table,
and relational attention over columns, rows, and PK--FK links lets the model
predict on unseen databases without updating its parameters
\citep{ranjan2026rt}. Later work improves pretraining and context: PluRel
synthesizes relational databases for pretraining
\citep{kothapalli2026plurel}, RT-J scales up pretraining for few-shot
prediction \citep{ranjan2026rtj}, and OpenRFM adds a batch-level in-context
learning layer to cope with sparse label coverage in sampled neighborhoods
\citep{chen2026openrfm}. KumoRFM-2 attends along the row, column, foreign-key,
and cross-sample dimensions \citep{hudovernik2026kumorfm2}. In \RT{}, Griffin,
and KumoRFM-2, attention over cells, rows, or examples determines which parts
of the relational context inform a prediction, and \RT{} keeps cell-level
tokens in every layer with about 22M parameters. \model{} instead lets the
database structure decide which records exchange information and learns only
how the exchanged messages are transformed, so it needs no attention and has
\mainParams{} parameters at its default width; it pools the
cells of each record once, giving up cell-level detail in exchange for
computation that scales with records rather than cells.

A second family flattens each entity's neighborhood into aggregate features
fixed in advance by deep feature synthesis and passes them to a tabular
in-context learner \citep{xu2026rdblearn,wang2026rdbpfn}. Their relational
encoding is therefore fixed rather than learned, and, like KumoRFM-2, they
rely on a support set of labeled
examples from the target task, whereas \RT{} and \model{} read only the
labels in the sampled context of each prediction; we follow the protocol of
\RT{}.
% Moved out of this paragraph (for an experiments-section comparison, if any):
% RDB-PFN reads 64-1,024 labeled examples and excludes the target entity's own
% label history; on our ten tasks its macro AUROC is 65.3 with 64 and 72.0 with
% 1,024 (wang2026rdbpfn Tables 6 and 10). RDBLearn reads up to 10,000.

\subsection{Transferable Relation Representations}
Knowledge graphs pose a related transfer problem, because the relation
vocabulary changes between training and inference graphs. InGram builds a
graph over relations whose edge weights reflect how many entities two
relations share, and derives embeddings of unseen relations from it
\citep{lee2023ingram}. ULTRA connects relations through four interactions
(head--head, tail--tail, head--tail, and tail--head), conditions the relation
representations on the query relation, and uses them for link prediction,
which enables zero-shot transfer across knowledge graphs
\citep{galkin2024ultra}. \model{} adapts this construction to relational
databases: a relation is a child--parent table pair, two relations are
adjacent when they share a table, and the resulting representations
condition message passing between records that carry their own content, for
entity classification rather than link prediction.

% ---- Previous version (before the 2026-11-02 rewrite), kept for reference ----
% \section{Related Work}
% \label{sec:related}
%
% \subsection{Relational Deep Learning}
% RDL turns a relational database into a graph of rows connected by PK--FK references and learns to aggregate information across tables \citep{fey2024rdl}.  RelBench standardizes databases, temporal prediction tasks, and evaluation splits for this setting \citep{robinson2024relbench}.  RelGNN uses atomic routes and composite message passing to better capture the structure of database graphs, including many-to-many relationships \citep{chen2025relgnn}. RelGT combines row features with type, hop, time, and local-structure information in a graph transformer \citep{dwivedi2026relgt}.  These approaches establish the value of explicitly modeling how rows are connected and how information passes between them. These models are typically trained for a specific database and prediction task, so applying them to a new database generally requires further training. Our work retains their row-graph representation and PK--FK message passing, but focuses on transfer to unseen databases without updating model parameters.
%
% \subsection{Foundation Models for Relational Data}
% Foundation models approach relational prediction through reusable representations and pretraining. This line includes both relational database foundation models and knowledge-graph foundation models; they share a transfer goal but differ in their inputs and prediction tasks.
%
% \subsection{Relational Database Foundation Models}
% Griffin combines shared encoders and decoders with relational message passing,
% but its downstream evaluation still requires task-specific fine-tuning
% \citep{wang2025griffin}. The Relational Transformer (\RT{}) removes this requirement by using task-table
% context and relational attention to predict on unseen databases without
% updating parameters \citep{ranjan2026rt}. Subsequent work addresses limitations
% in pretraining data and contextual labels. PluRel tackles the scarcity of
% diverse databases through synthetic relational data generation
% \citep{kothapalli2026plurel}. OpenRFM identifies sparse label coverage in
% \RT{}'s sampled neighborhoods and adds batch-level in-context learning
% \citep{chen2026openrfm}. RT-J improves few-shot prediction by scaling
% pretraining and revising context retrieval and training objectives
% \citep{ranjan2026rtj}.
%
% Alongside these improvements to pretraining and context, the granularity of
% relational computation remains a design choice. \RT{} maintains cell-level
% representations throughout its backbone. RDBLearn highlights the risk of
% losing task-relevant information through cross-column compression and
% instead preserves columns for a tabular in-context learner
% \citep{xu2026rdblearn}. We investigate whether pooling cells into compact row
% representations can still support competitive transfer when combined with
% explicit relation guidance. \model{} uses RelationGraph to guide RowGraph
% message passing and shared class-candidate nodes to aggregate labeled context
% for prediction.
% Griffin connects graph-centric RDL with foundation-model pretraining through shared encoders, relational message passing, and unified decoding \citep{wang2025griffin}.  It is an early point of contact between the two lines, with downstream transfer based on task adaptation.  \RT{} takes a frozen-transfer approach through task-table prompting, cell tokenization, and attention over columns, rows, and PK--FK neighborhoods \citep{ranjan2026rt}.  It is our closest comparator under the ten-task leave-database-out suite.
%
% Subsequent work explores several ways to improve reuse.  PluRel generates diverse synthetic schemas, PK--FK connectivity, and features for pretraining \citep{kothapalli2026plurel}.  RDBLearn combines a column-preserving relational encoder with a tabular in-context learner \citep{xu2026rdblearn}, while KumoRFM-2 combines row, column, FK, and cross-sample processing \citep{hudovernik2026kumorfm2}.  OpenRFM studies sparse contextual label coverage and introduces a batch-level support channel \citep{chen2026openrfm}; These methods demonstrate that relational foundation models encompass several representation choices. Our work investigates how relation-specific computation can be shared across database schemas, using structurally inferred relation states to guide prediction over compact row representations. We evaluate this design under the leave-database-out protocol used by \RT{}.
%
% \paragraph{Knowledge-graph foundation models.}
% Knowledge graphs expose a related transfer problem: entity and relation vocabularies can change between training and inference. Inductive methods such as InGram construct relation graphs to represent previously unseen relations \citep{lee2023ingram}.  Building on the goal of reusable graph reasoning, ULTRA is a knowledge-graph foundation model that constructs query-conditioned relation states through four endpoint interactions: head--head, tail--tail, head--tail, and tail--head.  These states condition entity-level link prediction and allow a pretrained model to transfer across graphs with different vocabularies \citep{galkin2024ultra}. This principle provides the foundation-model side of our relation encoder. Its application to databases requires heterogeneous cell encoding, time-respecting context construction, and an entity-classification readout in addition to structural reasoning.
%
% \paragraph{Connecting the two lines.}
% \model{} combines RDL's row-level representation with the transferable relation-construction principle of knowledge-graph foundation models and is evaluated against representative relational foundation models. We instantiate a relation node for each retained ordered child--parent table pair, initialize the relations instantiated by the query row's outgoing FKs, and use their inferred states to condition both directions of FK message passing. The relation identities are local to the input database, while the encoders, interaction operators, and readout are shared. The resulting model tests whether this combination can support competitive frozen transfer with a small row-graph backbone. As in relational in-context learning, frozen weights can still use historical labels supplied in the input; we therefore evaluate both Standard context and a controlled intervention on the query entity's target history.
